% !TEX root = ../main.tex
\begin{story}
本章を締めくくるにあたり, $g$-次ガルニエ系 $G_g$ と $2 \times 2$ シュレジンガー系との関係について調べておこう.

$$
\vec{y} = \begin{pmatrix} y_1 \\ y_2 \end{pmatrix}
$$

についての, シュレジンガー型方程式
\begin{equation} \tag{3.103}
    \frac{d\vec{y}}{dx} = \left( \frac{B_0}{x} + \frac{B_1}{x-1} + \sum_{l=1}^g \frac{A_l}{x-t_l} \right) \vec{y}
\end{equation}
を考える. このフックス型方程式について次のような仮定をおく.

$$
\det B_\Delta = \det A_l = 0 \qquad (\Delta = 0, 1, \quad l = 1, \cdots, g)
$$

\begin{equation} \tag{3.104}
    B_0 + B_1 + \sum_{l=1}^g A_l = - \begin{pmatrix} \chi_\infty & 0 \\ 0 & \chi_\infty + \kappa_\infty - 1 \end{pmatrix}
\end{equation}

前述のように, これは一般性を失わない. また

$$
\mathrm{Trace}\, B_\Delta = \kappa_\Delta, \qquad \mathrm{Trace}\, A_l = \theta_l
$$

とおく. これらのパラメータの間には, フックスの関係式

$$
\kappa_0 + \kappa_1 + \sum_{l=1}^g \theta_l + 2\chi_\infty + \kappa_\infty = 1
$$

が成り立っている.

ここで(3.103)のモノドロミー保存変形を考えると, シュレジンガー系
\end{story}